% GENERATED FILE. Edit canonical lessons, then run npm run generate:books.
% canonical-source-hash: fnv1a64:9ba61dfa7cc633c6
% canonical-lessons: JA-C431-tomadou, JA-C431-bikkuri-suru, JA-C431-akireru, JA-C431-odorokasu, JA-C431-harahara-suru, JA-R431-review-426-431

\chapter{Surprise and Suspense}
\label{ch:ja-surprise-and-suspense}
\hlchaptermodality{431}

\begin{chapteropening}
\textbf{By the end of this chapter:} \emph{I can say what startled or puzzled me.}

Complete the last lesson of chapter 431: I can say what startled or puzzled me.
\end{chapteropening}

\section[\texorpdfstring{\ja{とまどう} (tomadou)}{とまどう (tomadou)}]{\ja{とまどう} (tomadou) — to be bewildered; to be at a loss}

\label{lesson:JA-C431-tomadou}



\begin{quote}
Before the new one: say the Japanese for unpleasant, then the Japanese for to run wild.
\end{quote}

\subsection*{You'll want to know: \ja{とまどう}}
\textbf{\ja{とまどう}} — \emph{tomadou} — \textquotedblleft{}to be bewildered; to be at a loss\textquotedblright{}.

\begin{quote}
\textbf{\ja{せんせいの} \ja{しつもんに} \ja{とまどいました。}} — \emph{sensei no shitsumon ni tomadoimashita} — \textquotedblleft{}The teacher's question left me at a loss.\textquotedblright{}
\end{quote}

Polite \textbf{\ja{とまどいます}}, \emph{tomadoimasu}; a godan verb. Past (polite): \textbf{\ja{がくせいが} \ja{とまどいました。}}, \emph{gakusei ga tomadoimashita}.

\begin{cousinweb}
\textbf{\ja{とまどう}} (\emph{tomadou}) is a native Japanese word (wago), inherited rather than borrowed from Chinese: its second part is \emph{madou}, \textquotedblleft{}to be perplexed\textquotedblright{}, and JMdict spells the first, \emph{to}, with the kanji for \textquotedblleft{}door\textquotedblright{} or for \textquotedblleft{}route\textquotedblright{}, so which \emph{to} it is is not settled (JMdict, KANJIDIC2).
\end{cousinweb}

\begin{grammarlens}[title={What you've built}]
One more word for describing what you lived and felt.
\end{grammarlens}

\subsection*{Your turn}
\begin{itemize}
\raggedright
  \item \emph{Say it:} \emph{tomadou}
  \item \emph{Say it:} \emph{tomadou}, once more
  \item \emph{Say it:} \emph{sensei no shitsumon ni tomadoimashita} — The teacher's question left me at a loss.
  \item \emph{Recall:} say the Japanese for unpleasant, then the Japanese for to run wild, then say \emph{tomadou} again
\end{itemize}

\begin{tcolorbox}[breakable,colback=teal!4,colframe=teal!35!black,title={Before you move on}]
What does \ja{とまどう} mean? (\textquotedblleft{}To be bewildered; to be at a loss\textquotedblright{}.) Say it once more.
\end{tcolorbox}
\section[\texorpdfstring{\ja{びっくりする} (bikkuri suru)}{びっくりする (bikkuri suru)}]{\ja{びっくりする} (bikkuri suru) — to be startled}

\label{lesson:JA-C431-bikkuri-suru}



\begin{quote}
Before the new one: say the Japanese for to run wild, then the Japanese for to be bewildered.
\end{quote}

\subsection*{You'll want to know: \ja{びっくりする}}
\textbf{\ja{びっくりする}} — \emph{bikkuri suru} — \textquotedblleft{}to be startled\textquotedblright{}.

\begin{quote}
\textbf{\ja{かみなりの} \ja{おとに} \ja{びっくりしました。}} — \emph{kaminari no oto ni bikkuri shimashita} — \textquotedblleft{}I was startled by the sound of thunder.\textquotedblright{}
\end{quote}

Polite \textbf{\ja{びっくりします}}, \emph{bikkurishimasu}; a \ja{する} verb. Past (polite): \textbf{\ja{ねこが} \ja{びっくりしました。}}, \emph{neko ga bikkuri shimashita}.

\begin{cousinweb}
\textbf{\ja{びっくりする}} (\emph{bikkuri suru}) is a native Japanese mimetic word, \emph{bikkuri}, inherited rather than borrowed from Chinese, made a verb with \emph{suru}, \textquotedblleft{}to do\textquotedblright{}; its older origin is not settled (JMdict).
\end{cousinweb}

\begin{grammarlens}[title={What you've built}]
One more word for describing what you lived and felt.
\end{grammarlens}

\subsection*{Your turn}
\begin{itemize}
\raggedright
  \item \emph{Say it:} \emph{bikkuri suru}
  \item \emph{Say it:} \emph{bikkuri suru}, once more
  \item \emph{Say it:} \emph{kaminari no oto ni bikkuri shimashita} — I was startled by the sound of thunder.
  \item \emph{Recall:} say the Japanese for to run wild, then the Japanese for to be bewildered, then say \emph{bikkuri suru} again
\end{itemize}

\begin{tcolorbox}[breakable,colback=teal!4,colframe=teal!35!black,title={Before you move on}]
What does \ja{びっくりする} mean? (\textquotedblleft{}To be startled\textquotedblright{}.) Say it once more.
\end{tcolorbox}
\section[\texorpdfstring{\ja{あきれる} (akireru)}{あきれる (akireru)}]{\ja{あきれる} (akireru) — to be dumbfounded; to be appalled}

\label{lesson:JA-C431-akireru}



\begin{quote}
Before the new one: say the Japanese for to be bewildered, then the Japanese for to be startled.
\end{quote}

\subsection*{You'll want to know: \ja{あきれる}}
\textbf{\ja{あきれる}} — \emph{akireru} — \textquotedblleft{}to be dumbfounded; to be appalled\textquotedblright{}.

\begin{quote}
\textbf{\ja{かれの} \ja{うそに} \ja{みんな} \ja{あきれました。}} — \emph{kare no uso ni minna akiremashita} — \textquotedblleft{}Everyone was dumbfounded by his lie.\textquotedblright{}
\end{quote}

Polite \textbf{\ja{あきれます}}, \emph{akiremasu}; an ichidan verb. Past (polite): \textbf{\ja{せんせいが} \ja{あきれました。}}, \emph{sensei ga akiremashita}.

\begin{cousinweb}
\textbf{\ja{あきれる}} (\emph{akireru}) is a native Japanese word (wago), inherited rather than borrowed from Chinese; its kanji spelling was attached to it later (JMdict).
\end{cousinweb}

\begin{grammarlens}[title={What you've built}]
One more word for describing what you lived and felt.
\end{grammarlens}

\subsection*{Your turn}
\begin{itemize}
\raggedright
  \item \emph{Say it:} \emph{akireru}
  \item \emph{Say it:} \emph{akireru}, once more
  \item \emph{Say it:} \emph{kare no uso ni minna akiremashita} — Everyone was dumbfounded by his lie.
  \item \emph{Recall:} say the Japanese for to be bewildered, then the Japanese for to be startled, then say \emph{akireru} again
\end{itemize}

\begin{tcolorbox}[breakable,colback=teal!4,colframe=teal!35!black,title={Before you move on}]
What does \ja{あきれる} mean? (\textquotedblleft{}To be dumbfounded; to be appalled\textquotedblright{}.) Say it once more.
\end{tcolorbox}
\section[\texorpdfstring{\ja{おどろかす} (odorokasu)}{おどろかす (odorokasu)}]{\ja{おどろかす} (odorokasu) — to surprise someone; to startle}

\label{lesson:JA-C431-odorokasu}



\begin{quote}
Before the new one: say the Japanese for to be startled, then the Japanese for to be dumbfounded.
\end{quote}

\subsection*{You'll want to know: \ja{おどろかす}}
\textbf{\ja{おどろかす}} — \emph{odorokasu} — \textquotedblleft{}to surprise someone; to startle\textquotedblright{}.

\begin{quote}
\textbf{\ja{おとうとが} \ja{うしろから} \ja{ははを} \ja{おどろかしました。}} — \emph{otōto ga ushiro kara haha o odorokashimashita} — \textquotedblleft{}My younger brother startled my mother from behind.\textquotedblright{}
\end{quote}

Polite \textbf{\ja{おどろかします}}, \emph{odorokashimasu}; a godan verb. Past (polite): \textbf{\ja{ねこが} \ja{こどもを} \ja{おどろかしました。}}, \emph{neko ga kodomo o odorokashimashita}.

\begin{cousinweb}
\textbf{\ja{おどろかす}} (\emph{odorokasu}) is a native Japanese word (wago), inherited rather than borrowed from Chinese: the transitive partner of \emph{odoroku}, \textquotedblleft{}to be surprised\textquotedblright{}, and its kanji was attached later as a kun'yomi (JMdict, KANJIDIC2).
\end{cousinweb}

\begin{grammarlens}[title={What you've built}]
One more word for describing what you lived and felt.
\end{grammarlens}

\subsection*{Your turn}
\begin{itemize}
\raggedright
  \item \emph{Say it:} \emph{odorokasu}
  \item \emph{Say it:} \emph{odorokasu}, once more
  \item \emph{Say it:} \emph{otōto ga ushiro kara haha o odorokashimashita} — My younger brother startled my mother from behind.
  \item \emph{Recall:} say the Japanese for to be startled, then the Japanese for to be dumbfounded, then say \emph{odorokasu} again
\end{itemize}

\begin{tcolorbox}[breakable,colback=teal!4,colframe=teal!35!black,title={Before you move on}]
What does \ja{おどろかす} mean? (\textquotedblleft{}To surprise someone; to startle\textquotedblright{}.) Say it once more.
\end{tcolorbox}
\section[\texorpdfstring{\ja{はらはらする} (harahara suru)}{はらはらする (harahara suru)}]{\ja{はらはらする} (harahara suru) — to be on edge; to be in suspense}

\label{lesson:JA-C431-harahara-suru}



\begin{quote}
Before the new one: say the Japanese for to be dumbfounded, then the Japanese for to surprise someone.
\end{quote}

\subsection*{You'll want to know: \ja{はらはらする}}
\textbf{\ja{はらはらする}} — \emph{harahara suru} — \textquotedblleft{}to be on edge; to be in suspense\textquotedblright{}.

\begin{quote}
\textbf{\ja{おとうとの} \ja{しあいで} \ja{はらはらしました。}} — \emph{otōto no shiai de harahara shimashita} — \textquotedblleft{}I was on edge at my younger brother's match.\textquotedblright{}
\end{quote}

Polite \textbf{\ja{はらはらします}}, \emph{haraharashimasu}; a \ja{する} verb. Past (polite): \textbf{\ja{ははが} \ja{はらはらしました。}}, \emph{haha ga harahara shimashita}.

\begin{cousinweb}
\textbf{\ja{はらはらする}} (\emph{harahara suru}) is a native Japanese mimetic word, \emph{harahara}, inherited rather than borrowed from Chinese, made a verb with \emph{suru}, \textquotedblleft{}to do\textquotedblright{}; its older origin is not settled (JMdict).
\end{cousinweb}

\begin{grammarlens}[title={What you've built}]
One more word for describing what you lived and felt.
\end{grammarlens}

\subsection*{Your turn}
\begin{itemize}
\raggedright
  \item \emph{Say it:} \emph{harahara suru}
  \item \emph{Say it:} \emph{harahara suru}, once more
  \item \emph{Say it:} \emph{otōto no shiai de harahara shimashita} — I was on edge at my younger brother's match.
  \item \emph{Recall:} say the Japanese for to be dumbfounded, then the Japanese for to surprise someone, then say \emph{harahara suru} again
\end{itemize}

\begin{tcolorbox}[breakable,colback=teal!4,colframe=teal!35!black,title={Before you move on}]
What does \ja{はらはらする} mean? (\textquotedblleft{}To be on edge; to be in suspense\textquotedblright{}.) Say it once more.
\end{tcolorbox}
\section[(dialogue)]{Looking back: from How I Feel Today to Surprise and Suspense}

\label{lesson:JA-R431-review-426-431}



\begin{quote}
Say the Japanese for one's mood and anxious.
\end{quote}

\subsection*{Your turn}
Say these aloud:

\begin{itemize}
\raggedright
  \item one's mood, anxious, unbothered, refreshed, to endure
  \item to get nervous, to have a pounding heart, to be restless, to get irritated, to hesitate
  \item to be thrilled, to get worked up, to be deeply moved, to frolic, to be full of energy
  \item to be disappointed, to feel down, a sigh, hard to bear, to comfort
  \item anger, to curse at, to glare at, unpleasant, to run wild
  \item to be bewildered, to be startled, to be dumbfounded, to surprise someone, to be on edge
\end{itemize}

\begin{tcolorbox}[breakable,colback=teal!4,colframe=teal!35!black,title={Before you move on}]
One's mood? (\textbf{\ja{きぶん}}, \emph{kibun}.) Anxious? (\textbf{\ja{ふあん}}, \emph{fuan}.) Unbothered? (\textbf{\ja{へいき}}, \emph{heiki}.) Refreshed? (\textbf{\ja{すっきり}}, \emph{sukkiri}.) To endure? (\textbf{\ja{たえる}}, \emph{taeru}.) To get nervous? (\textbf{\ja{きんちょうする}}, \emph{kinchō suru}.) To have a pounding heart? (\textbf{\ja{どきどきする}}, \emph{dokidoki suru}.) To be restless? (\textbf{\ja{そわそわする}}, \emph{sowasowa suru}.) To get irritated? (\textbf{\ja{いらいらする}}, \emph{iraira suru}.) To hesitate? (\textbf{\ja{ためらう}}, \emph{tamerau}.) To be thrilled? (\textbf{\ja{わくわくする}}, \emph{wakuwaku suru}.) To get worked up? (\textbf{\ja{こうふんする}}, \emph{kōfun suru}.) To be deeply moved? (\textbf{\ja{かんどうする}}, \emph{kandō suru}.) To frolic? (\textbf{\ja{はしゃぐ}}, \emph{hashagu}.) To be full of energy? (\textbf{\ja{はりきる}}, \emph{harikiru}.) To be disappointed? (\textbf{\ja{がっかりする}}, \emph{gakkari suru}.) To feel down? (\textbf{\ja{おちこむ}}, \emph{ochikomu}.) A sigh? (\textbf{\ja{ためいき}}, \emph{tameiki}.) Hard to bear? (\textbf{\ja{つらい}}, \emph{tsurai}.) To comfort? (\textbf{\ja{なぐさめる}}, \emph{nagusameru}.) Anger? (\textbf{\ja{いかり}}, \emph{ikari}.) To curse at? (\textbf{\ja{ののしる}}, \emph{nonoshiru}.) To glare at? (\textbf{\ja{にらむ}}, \emph{niramu}.) Unpleasant? (\textbf{\ja{いや}}, \emph{iya}.) To run wild? (\textbf{\ja{あばれる}}, \emph{abareru}.) To be bewildered? (\textbf{\ja{とまどう}}, \emph{tomadou}.) To be startled? (\textbf{\ja{びっくりする}}, \emph{bikkuri suru}.) To be dumbfounded? (\textbf{\ja{あきれる}}, \emph{akireru}.) To surprise someone? (\textbf{\ja{おどろかす}}, \emph{odorokasu}.) To be on edge? (\textbf{\ja{はらはらする}}, \emph{harahara suru}.)
\end{tcolorbox}
